% Minimal preamble for P2 cards taken from sources/comparability-partitions-2408.13523v2/source0.tex
% (arXiv:2408.13523v2, "Partitioning perfect graphs into comparability graphs").
% Every line below is copied mechanically from that file; nothing was invented, renamed or redefined.
%   line 2    : \documentclass[12pt]{article}
%   line 3    : \usepackage{amssymb,latexsym}
%   line 4    : \usepackage{amsmath}
%   line 5    : \usepackage[dvipdfm]{graphicx}  -- the dvipdfm driver option is dropped, because the
%               assembler forces the xelatex engine and a forced driver option contradicts it.
%   line 16   : \usepackage{amsfonts}
%   line 17   : \usepackage{amsthm}
%   lines 20-30: the author's theorem declarations
%   lines 32-41: the author's own macro definitions
% Omitted source preamble lines and why they are not needed for these body spans:
%   line 1 (provenance comment), lines 6-11 (tikz and its libraries: no drawing in the retained spans),
%   lines 7 and 9 (verbatim), lines 12-13 (page geometry), line 14 (rawfonts),
%   line 15 (psfrag), line 18 (hyperref with the pagebackref option: the retained spans print no
%   citation anchors and the assembly adds its own hyperref machinery), lines 43-61 (the author's
%   backref neutralisation and his margin-note macros, which depend on the omitted hyperref).
% The retained spans of the body use only \em, \bf, \cal, amsthm environments, amsmath displays and
% the macros defined in lines 32-41 below.
\documentclass[12pt]{article}
\usepackage{amssymb,latexsym}
\usepackage{amsmath}
\usepackage{graphicx}
\usepackage{amsfonts}
\usepackage{amsthm}

\newtheorem{theorem}{Theorem}%[section]
\newtheorem{problem}[theorem]{Problem}%[section]
\newtheorem{observation}[theorem]{Observation}
\newtheorem{lemma}[theorem]{Lemma}
\newtheorem{corollary}[theorem]{Corollary}
\newtheorem{definition}{Definition}
\newtheorem{example}[theorem]{Example}
\newtheorem{conjecture}{Conjecture}%[section]
\newtheorem{proposition}[theorem]{Proposition}
\newtheorem{claim}[theorem]{Claim}
\newtheorem{remark}[theorem]{Remark}

\def\qed{\ifhmode\unskip\nobreak\fi\quad\ifmmode\Box\else$\Box$\fi}
\def\pont{\hspace{-6pt}{\bf.\ }}
\def\qed{\hfill $\Box$}

\def\R{{\mbox{\boldmath$\mathcal R$}}}
% the author's second definition of \R, left commented out exactly as in the source
\def\F{{\mathcal F}}
\def\P{{\mathcal P}}

\def\newpage{\vfill\eject}
\title{Partitioning a perfect graph into bipartite graphs and comparability graphs from its clique and stability numbers}
\author{András Gyárfás, Márton Marits, Géza Tóth}
\date{arxiv:2408.13523v2; CC BY 4.0}
\begin{document}
\maketitle
\noindent\textit{Editorial note.} The counted unit is the article's own Corollary genupbound at source lines 159-164 together with the author's proof that opens with the bold header Proof. at source line 166 and closes with the square at source line 175; the author separates the two by the sentence at source line 158, which introduces the corollary, so the pairing is by adjacency through the intervening blank line only.\par\medskip
\noindent\textit{Editorial note.} The retained context at source lines 100-105 is the article's own definition of perfect graphs and of the four parameters omega, chi, alpha and theta; the retained context at source lines 117-120 and 128-129 is its own description of comparability graphs and its definitions of p(G) and c(G).\par\medskip
\noindent\textit{Editorial note.} The retained context at source lines 154-155 is the article's own statement of Theorem A on partitioning an edge set into bipartite graphs, which the proof applies twice. The surrounding sentences that carry bibliographic citations are deliberately not retained, so the material carries no bibliography.\par\medskip
\noindent\textit{Editorial note.} The source preamble cannot be retained as it stands: it loads rawfonts, psfrag, verbatim, tikz with several libraries, hyperref with the pagebackref option and a dvi driver option for graphicx, none of which the retained body spans use, and it redefines backref commands that only exist because of that hyperref option. The retained preamble file therefore carries, line for line, source line 2, source line 3, source line 4, source line 16, source line 17, the theorem declarations of source lines 20-30 and the author macro definitions of source lines 32-41; the only textual change anywhere is the removal of the dvipdfm driver option from the graphicx line, because the assembly forces the xelatex engine. No new macro, package or declaration was introduced.\par\medskip
\noindent\textit{Editorial note.} The counted corollary's own itemize environment is standard LaTeX and needs no additional package. The article's other results, including Theorem twocover, Theorem gensplit, Theorem bipline, Theorem unit, Theorem co-triangulated, Theorem loglog, Theorem ub, Lemma rec and Proposition H9, are not counted here.\par\medskip
\noindent\textit{Allowed standard prerequisites (editorial).}\begin{itemize}
\item Perfect graphs and the parameters omega, chi, alpha and theta. A graph is perfect when every induced subgraph H satisfies omega(H)=chi(H), where omega is the size of a largest complete subgraph and chi is the chromatic number; equivalently the complement satisfies alpha(H)=theta(H), where alpha is the size of a largest independent set and theta is the minimum number of cliques covering the vertex set; retained at source lines 100-105.. The counted corollary is stated for perfect graphs and bounds the two partitions by the ceiling of the logarithm of omega(G) and of alpha(G), so all four parameters and their equality must be available.
\item Theorem A on partitioning an edge set into bipartite graphs. The minimum number of bipartite graphs needed to partition the edge set of a graph G is the ceiling of the base-2 logarithm of chi(G); the article states this as Theorem A at source lines 154-155 and attributes it to Harary, Hsu and Miller.. The proof applies this theorem twice: once to G itself using chi(G)=omega(G), and once to the graph H left after deleting the edges of a clique cover, using chi(H) <= alpha(G).
\item Weak perfect graph theorem and the clique cover by alpha(G) complete subgraphs. For a perfect graph the complement is perfect, so the vertex set can be partitioned into alpha(G) complete subgraphs, and the union of those complete subgraphs is a single comparability graph; the article invokes this in the proof at source lines 170-172 and describes comparability graphs at source lines 117-120.. Gives the one comparability graph of part (ii) and the estimate chi(H) <= alpha(G) for the graph H obtained by deleting its edges.
\end{itemize}\medskip
Perfect graphs are the graphs $G$ in which all induced subgraphs $H\subseteq G$
satisfy the property $\omega(H)=\chi(H)$ where $\omega(H)$ and $\chi(H)$ denote
the size of the largest complete subgraph and the chromatic number.
Berge put forward two conjectures, the weaker one was that complements of
perfect graphs are also perfect. This means that the induced subgraphs $H$ also satisfy $\alpha(H)=\theta(H)$,
where $\alpha(H)$ is the cardinality of the largest independent set and $\theta(H)$ is the minimum number of cliques covering $V(H)$.

Here we focus on a well known family of perfect graphs, the {\em comparability graphs},
which can be obtained from a partially ordered set as follows.
Vertices correspond to the elements of the partially ordered set,
and two vertices are connected by an undirected edge if and only if the corresponding elements

Let $p(G)$ denote the minimum number $m$ such that $E(G)=\cup_{i=1}^m E(G_i)$ where the $G_i$-s are edge-disjoint comparability subgraphs of $G$.
Dropping  the condition that the $G_i$-s are edge-disjoint, we define $c(G)$ as the smallest number of comparability subgraphs of $G$ needed to {\em cover} $E(G)$.

    {\bf Theorem A.} The minimum number of bipartite graphs needed to partition the edge set of a graph $G$ is
    $\lceil\log(\chi(G))\rceil$, where the logarithm is base 2 (we use this through the paper).

\begin{corollary}\label{genupbound}  Every perfect graph $G$ can be partitioned into
\begin{itemize}
\item (i) $\lceil\log(\omega(G))\rceil$ bipartite graphs
\item (ii) $1+\lceil\log(\alpha(G))\rceil$ comparability graphs
\end{itemize}
\end{corollary}

\noindent{\bf Proof. }
Since  $\chi(G)=\omega(G)$
for any perfect graph $G$, Theorem A implies that we can partition $E(G)$ into
$\lceil\log(\omega(G))\rceil$ bipartite graphs, proving (i).
Also, (by the weak perfect graph theorem) $V(G)$ can be partitioned
into $\alpha(G)$ complete subgraphs, whose union is {\em one}
comparability graph. Deleting the edges of these complete subgraphs,
we get a graph $H$ with $\chi(H)\le \alpha(G)$. Applying Theorem A
for $H$, we get at most $\lceil\log(\alpha(G))\rceil$ further
(bipartite) comparability graphs, proving (ii). \qed

\end{document}
